\documentclass[12pt,a4paper,twoside]{article}
\usepackage{bcc2003,url,graphicx,color,paralist,calc,rotating,psfrag,overpic}

\bcctitle{Counting Lattice Triangulations}
%\bccshorttitle{Lattice Triangulations}
\bccname{Volker Kaibel and G\"unter~M. Ziegler}
\bccaddress{Institut f\"ur Mathematik, MA 6--2\\
Technische Universit\"at Berlin\\
%Stra\ss e des 17. Juni 136\\
D--10623~Berlin\\
Germany}
\bccemail{\{kaibel,ziegler\}@math.tu-berlin.de}

% add new thoerem-like environments here
% \newtheorem{hyp}[theorem]{Hypotheis}
% \plaintheorems

% why not use \mathbb{R} etc? - to be discussed
\newcommand{\R}{\mathbb{R}}%{{\bf R}}
\newcommand{\Z}{\mathbb{Z}}%{{\bf Z}}
\newcommand{\N}{\mathbb{N}}%{{\bf N}}
\newcommand{\order}[1]{\mbox{O}({#1})}
\newcommand{\orderBig}[1]{\mbox{O}\Big({#1}\Big)}
\DeclareMathOperator{\littleo}{o}
\newcommand{\conv}{\mbox{\rm conv}}
\newcommand{\aff}{\mbox{\rm aff}}
\newcommand{\freg}{f^{\mathrm{reg}}}
\newcommand{\firreg}{f^{\mathrm{irreg}}}
\newcommand{\creg}{c^{\mathrm{reg}}}
\newcommand{\abov}{\prec}
\newcommand{\cplex}{\textsf{CPLEX}~\texttt{6.6.1}}
\newcommand{\topcom}{\textsf{TOPCOM}}

\newmathop{\deven}{{\rm DE}}
\newmathop{\dodd}{{\rm DO}}
\newmathop{\aut}{{\rm Aut}}
\newmathbin{\wreath}{{\rm wr}}


\begin{document}
\makebcctitle

\begin{abstract}
We discuss the problem
to count, or, more modestly, to estimate the number $f(m,n)$
of unimodular triangulations of the planar grid of size~$m\times n$.

Among other tools, we employ recursions
that allow one to compute the (huge) number of 
triangulations for small~$m$ and rather large~$n$ by
dynamic programming;
we show that this computation can be done in polynomial
time if $m$ is fixed, and present computational results
from our implementation of this approach.

We also present new upper and
lower bounds for large $m$ and~$n$, and we 
report about results obtained from a computer simulation of the random
walk that is generated by flips.
\end{abstract}

\input intro

\input values

\input bounds

\input trias

\subsection*{Acknowledgements.}
We are grateful to 
Oswin Aichholzer,
\'Emile Anclin, 
Stephan Meyer, and
J\"org Rambau
for valuable discussions. Thanks to Francisco Santos and to the
referee for helpful comments
on an earlier version of the manuscript.

\input biblio


\myaddress
\newpage

\input appendix

\end{document}

